\section{Discussion}
\label{sec:discussion}

Three studies have been reported: an observational comparison, a controlled
injection experiment, and a human preference study. Together they distinguish
targeted affected-code grounding from general review quality.

\subsection{Reconciling Capability and Aggregate Gains}
\label{sec:discussion-reconciliation}

The two experiments appear to disagree about effect size by an order of
magnitude. They estimate different quantities. Experiment~2 conditions on
the mechanism applying: every injected defect was placed so that its
consequence lay outside the edited file. It measures a capability (can the
reviewer see such a consequence at all?) and answers that a diff-only
reviewer cannot and a graph-equipped one often can (18/28 in this injection
sample). Experiment~1 samples pull requests without conditioning on that
mechanism, and its average graph effect is correspondingly smaller.

The graph gives the reviewer a new ability, while its average effect is small
on the 35 pull requests in Experiment~1. The experiment does not estimate how
often that ability would matter in the broader population of pull requests.

The join is conceptual rather than based on a shared criterion label.
Experiment~1's advantage concentrates in the nine structurally relevant
criteria, with integration (F3) and breaking-change impact (F4) among the
largest movers. Experiment~2 directly tests whether a review names an affected
dependent outside the diff. The human study's primary item (H1) asks which
review better names affected code.
These results converge on affected-code grounding. The exploratory clarity
item (H5) favours the baseline, which
further limits any claim of general review improvement.

\subsection{Complementarity of Structure and Similarity}
\label{sec:discussion-complementarity}

Retrieval answers a similarity question: which parts of this repository
resemble the change. That lifts many criteria a little. A dependency graph
answers a reachability question: what is connected to what changed. That
helps a narrow set of criteria a great deal and most criteria not at all.

Experiment~2 shows the two are not interchangeable. Retrieval scored 5/28
on cross-file defects against the graph's 18/28. A file that breaks because
it calls a renamed function need not resemble the file that changed.
Retrieval is not a substitute for structural cross-file detection, although it
remains competitive on broader rubric coverage.

The hybrid arm detects twenty-one of twenty-eight, numerically above both
inputs. The paired difference against the graph is not significant
($p = 0.375$), so it provides evidence of neither superiority nor dilution.
On the pull requests the picture is less favourable: the hybrid trails the
leading single arm on both scales and both arms on the total
(Table~\ref{tab:exp1-deltas}), so under naive concatenation the rubric
shows no additive gain.
Together with the crossing pattern in Experiment~1, the result is consistent
with complementary mechanisms, but does not show that naive concatenation
combines their benefits.

\subsection{Implications for Evaluation Design}
\label{sec:discussion-measurement}

The rubric estimate has two measured weaknesses. It is panel-dependent: the
same reviews yield effects ranging from significant to not, depending on
which judges are polled. And it is generator-dependent: three substitute
generators yield no significant effect
(Appendix~\ref{app:supporting-robustness}).

The practical recommendation generalises past this thesis: a
generate-then-LLM-judge design should not be the sole evidence for a context
mechanism. Where a mechanism can be stated precisely enough to build an
oracle, the oracle should be built.

\subsection{Practical Implications}
\label{sec:discussion-practical}

Commercial review tools such as CodeRabbit and Greptile already index
repositories and hand the model pre-digested dependency
findings~\cite{coderabbit2026,greptile2026}; this thesis shows that a model
given raw graph context can itself name affected code outside the diff, with
the gain localised to cross-file criteria. Three positions follow for
pipeline builders.

\emph{Supply the graph selectively.} The graph earns its place on changes
with cross-file consequences and contributes little elsewhere. A pull
request's structural fan-out can be computed before any review is generated,
so the decision whether to include graph context can itself be automated.

\emph{Resolve structure at function granularity where the analyser supports
it.} The human study's graph block carries file-level import edges, and a
module-level import does not imply that a change reaches the importer. This is a precision argument, not evidence that
function-level edges are always necessary: in Experiment~2 the file-level
dependency list alone recovers sixteen of the graph arm's eighteen
detections.

\emph{Invest in the analyser, not only the model.} Within Experiment~2 the
detection ladder moves substantially when the analyser changes and the model
is held fixed. Whether analyser improvements should take priority for
broader pull-request review remains untested.
